
\documentclass{article}
\usepackage{colortbl}
\usepackage{makecell}
\usepackage{multirow}
\usepackage{supertabular}

\begin{document}

\newcounter{utterance}

\centering \large Interaction Transcript for game `cladder', experiment `full\_v1.5\_default', episode 8317 with refineonly{-}llama{-}3.1{-}8b{-}instruct.
\vspace{24pt}

{ \footnotesize  \setcounter{utterance}{1}
\setlength{\tabcolsep}{0pt}
\begin{supertabular}{c@{$\;$}|p{.15\linewidth}@{}p{.15\linewidth}p{.15\linewidth}p{.15\linewidth}p{.15\linewidth}p{.15\linewidth}}
    \# & \multicolumn{2}{c}{Player} && \multicolumn{2}{c}{Game Master} \\
    \hline

    \theutterance \stepcounter{utterance}  
    & & & \multicolumn{4}{p{0.6\linewidth}}{
        \cellcolor[rgb]{0.9,0.9,0.9}{
            \makecell[{{p{\linewidth}}}]{
                \texttt{\tiny{[P1$\langle$GM]}}
                \texttt{You are an expert in causal inference. The following question is not a typical commonsense query, but rather a meticulously designed question created by a professor specializing in causal inference, intended to assess the students' mastery of the course content.} \\
\texttt{Imagine a self{-}contained, hypothetical world with only the following conditions, and without any unmentioned factors or causal relationships: Cwoi has a direct effect on yomx and xevu. Yomx has a direct effect on gwet. Gwet has a direct effect on xevu. Cwoi is unobserved. For those who are not yomx, the probability of gwet is 70\%. For those who are yomx, the probability of gwet is 44\%. For those who are not yomx and are not gwet, the probability of xevu is 64\%. For those who are not yomx and are gwet, the probability of xevu is 49\%. For those who are yomx and are not gwet, the probability of xevu is 58\%. For those who are yomx and are gwet, the probability of xevu is 89\%. The overall probability of yomx is 43\%. Does yomx negatively affect xevu through gwet?} \\
\texttt{Start your answer with 'Yes' or 'No', followed by additional reasoning or evidence to support your explanation.} \\
            }
        }
    }
     \\ \\

    \theutterance \stepcounter{utterance}  
    & & \multicolumn{4}{p{0.6\linewidth}}{
        \cellcolor[rgb]{0.95,0.95,0.95}{
            \makecell[{{p{\linewidth}}}]{
                \texttt{\tiny{[P1$\leftrightarrow$P]}}
                \texttt{{[}DRAFT{]}} \\
\texttt{No, yomx does not negatively affect xevu through gwet.} \\
\\ 
\texttt{To assess the causal relationship between yomx and xevu through gwet, we can use the backdoor criterion, which involves checking for confounding and mediation. } \\
\\ 
\texttt{First, we need to establish the causal relationship between yomx and gwet. The probability of gwet is 70\% for those who are not yomx and 44\% for those who are yomx. This indicates a negative association between yomx and gwet. } \\
\\ 
\texttt{Next, we need to establish the causal relationship between gwet and xevu. The probability of xevu is 64\% for those who are not yomx and are not gwet, 49\% for those who are not yomx and are gwet, 58\% for those who are yomx and are not gwet, and 89\% for those who are yomx and are gwet. This indicates a positive association between gwet and xevu.} \\
\\ 
\texttt{However, to determine if yomx negatively affects xevu through gwet, we need to examine the effect of yomx on the probability of xevu, conditional on gwet. } \\
\\ 
\texttt{For those who are not gwet, the probability of xevu is 64\% for those who are not yomx and 58\% for those who are yomx. This indicates a positive association between yomx and xevu, conditional on not being gwet.} \\
\\ 
\texttt{For those who are gwet, the probability of xevu is 49\% for those who are not yomx and 89\% for those who are yomx. This indicates a positive association between yomx and xevu, conditional on being gwet.} \\
\\ 
\texttt{Since the association between yomx and xevu is positive, both conditional on being gwet and not being gwet, we can conclude that yomx does not negatively affect xevu through gwet.} \\
            }
        }
    }
    & \\ \\

    \theutterance \stepcounter{utterance}  
    & \multicolumn{4}{p{0.6\linewidth}}{
        \cellcolor[rgb]{0.561,0.737,0.561}{
            \makecell[{{p{\linewidth}}}]{
                \texttt{\tiny{[P1$\rangle$GM]}}
                \texttt{No, yomx does not negatively affect xevu through gwet.} \\
\\ 
\texttt{To assess the causal relationship between yomx and xevu through gwet, we can use the backdoor criterion, which involves checking for confounding and mediation. } \\
\\ 
\texttt{First, we need to establish the causal relationship between yomx and gwet. The probability of gwet is 70\% for those who are not yomx and 44\% for those who are yomx. This indicates a negative association between yomx and gwet.} \\
\\ 
\texttt{Next, we need to establish the causal relationship between gwet and xevu. The probability of xevu is 64\% for those who are not yomx and are not gwet, 49\% for those who are not yomx and are gwet, 58\% for those who are yomx and are not gwet, and 89\% for those who are yomx and are gwet. This indicates a positive association between gwet and xevu.} \\
\\ 
\texttt{However, to determine if yomx negatively affects xevu through gwet, we need to examine the effect of yomx on the probability of xevu, conditional on gwet. } \\
\\ 
\texttt{For those who are not gwet, the probability of xevu is 64\% for those who are not yomx and 58\% for those who are yomx. This indicates a positive association between yomx and xevu, conditional on not being gwet.} \\
\\ 
\texttt{For those who are gwet, the probability of xevu is 49\% for those who are not yomx and 89\% for those who are yomx. This indicates a positive association between yomx and xevu, conditional on being gwet.} \\
\\ 
\texttt{Since the association between yomx and xevu is positive, both conditional on being gwet and not being gwet, we can conclude that yomx does not negatively affect xevu through gwet.} \\
            }
        }
    }
    & & \\ \\

    \theutterance \stepcounter{utterance}  
    & & & \multicolumn{2}{p{0.3\linewidth}}{
        \cellcolor[rgb]{0.9,0.9,0.9}{
            \makecell[{{p{\linewidth}}}]{
                \texttt{\tiny{[GM$|$GM]}}
                \texttt{No} \\
            }
        }
    }
    & & \\ \\

    \theutterance \stepcounter{utterance}  
    & & & \multicolumn{2}{p{0.3\linewidth}}{
        \cellcolor[rgb]{0.9,0.9,0.9}{
            \makecell[{{p{\linewidth}}}]{
                \texttt{\tiny{[GM$|$GM]}}
                \texttt{yes} \\
            }
        }
    }
    & & \\ \\

    \theutterance \stepcounter{utterance}  
    & & & \multicolumn{2}{p{0.3\linewidth}}{
        \cellcolor[rgb]{0.9,0.9,0.9}{
            \makecell[{{p{\linewidth}}}]{
                \texttt{\tiny{[GM$|$GM]}}
                \texttt{game\_result = LOSE} \\
            }
        }
    }
    & & \\ \\

\end{supertabular}
}

\end{document}
